\documentclass{article}
\usepackage{amsmath,amssymb,amsthm}
\begin{document}

\textbf{Subproblem 3 --- Long chains force extra surviving points, hence $n\le 4$.}

\medskip
\textbf{Context.}
Use the same setup and equality-case assumptions as Subproblems~1--2, whose
results may be freely used.  In particular: every surviving point is lonely
($F(\ell_s)=\varepsilon(t_s)\in\{\pm1\}$), the boundary of $\operatorname{conv}(S)$
splits into an upper chain $U$ and a lower chain $L$, and every vertex of
$\operatorname{conv}(S)+\operatorname{conv}(T)$ lies in the support of $F$.

\medskip
\textbf{Statement.}  Let $C\in\{U,L\}$ be one of the two boundary chains with
$m = |C|$ vertices ($m\ge 2$ since $C$ shares its two endpoints with the other
chain).  Prove: if $m\ge 4$, then the number of surviving points
$|\{p:F(p)\neq 0\}|$ is strictly greater than $n$, giving a contradiction.
Hence both $|U|\le 3$ and $|L|\le 3$, and therefore $n = |U|+|L|-2 \le 4$.

\medskip
\textbf{Key idea.}
On each strictly convex arc $C$, the outer unit-normal direction sweeps
monotonically.  As the normal direction rotates from $s_a$-end to $s_b$-end
along $C$, the exposed vertex of $\operatorname{conv}(T)$ also traces a monotone
path among the vertices of $\operatorname{conv}(T)$.  If $C$ has sign changes
(which it does, by Subproblem~2), these sign changes force new \emph{vertices}
of $P+Q$ that are not among the $n$ lonely points, contradicting the equality.
Specifically: each \emph{sign change} along the chain (from $+$ to $-$ or
$-$ to $+$) combined with the strict convexity of the chain and the monotone
motion of the exposed $T$-vertex forces at least one extra vertex of $P+Q$
to appear between the two adjacent chain vertices.  A chain of length $m\ge 4$
has at least one interior sign change, hence at least one such extra vertex,
giving $|\operatorname{supp}(F)|\ge n+1$.

\medskip
\textbf{Hint.}
The vertices of $P+Q$ are in bijection with the \emph{edge-normal directions} of
$P+Q$, which are all normal directions appearing in $P$ or $Q$.  Between two
consecutive vertices $v_i, v_{i+1}$ of $C$ with $\varepsilon(t_{v_i})\neq
\varepsilon(t_{v_{i+1}})$ (a sign change), the corresponding exposed
translate in $T$ must change: there is a direction in the normal-cone interval
$[n_i, n_{i+1}]$ (the normal directions of the edge $v_iv_{i+1}$ of $P$) at
which the exposed vertex of $Q$ transitions.  This transition produces a new
vertex of $P+Q$ in the interior of the edge $v_iv_{i+1}+(t_{v_i},t_{v_{i+1}})$
that is not one of the $n$ lonely points.  Since vertices of $P+Q$ always have
$F\neq 0$ (Subproblem~1, item~6), this gives an extra nonzero point.

\end{document}
